\documentclass[11pt]{article}

\usepackage[margin=1.1in]{geometry}
\usepackage{amsmath,amssymb,amsthm,mathrsfs}
\usepackage[colorlinks=true,linkcolor=black,urlcolor=black,citecolor=black]{hyperref}

\newtheorem{theorem}{Theorem}
\newtheorem{lemma}[theorem]{Lemma}
\newtheorem{proposition}[theorem]{Proposition}
\newtheorem{corollary}[theorem]{Corollary}
\theoremstyle{definition}
\newtheorem{definition}[theorem]{Definition}
\theoremstyle{remark}
\newtheorem{remark}[theorem]{Remark}

\newcommand{\C}{\mathbf C}
\newcommand{\Q}{\mathbf Q}
\newcommand{\Z}{\mathbf Z}
\newcommand{\PP}{\mathbf P}

\title{Frame transport for framed formal monodromy:\\
  the lemma, its hypotheses, and the receiver problem\\
  \large Fifth version, closing the coalesced case}
\author{Memo for a reader of \emph{Irrationality of Cubic Threefolds after One
    Stabilization}}
\date{15 August 2026}

\begin{document}
\title{Frame transport memo --- archive}
\author{Tavis Rudd}
\date{August 2026}
\maketitle

\noindent Append-only history for
\texttt{2026-08-15-c912-frame-transport-memo.tex}.  Superseded routes,
closed checks, and correction trails are kept here verbatim so the live
memo can state current conclusions only.  Cross-references to statements
that remain in the live memo are rendered as bracketed label names.


\section{The scissors: hypotheses with no common instance}
\noindent\emph{Archived from the live memo, lines 392--421.}

\subsection*{The scissors: these hypotheses have no common instance in the
application}


Theorem~\textup{[thm:transport]} is true, and in the application it is empty. Its two
hypotheses pull in opposite directions.

\begin{itemize}
\item If $\Gamma=0$, so $\mathcal H=\Omega_V(\!(z)\!)$, then Levelt--Turrittin
  supplies solution matrices, but the pro-Laurent comparison gauge has
  unbounded negative $z$-order and is not in $\mathrm{GL}_n(\mathcal H)$.
\item If $\Gamma\ne0$ with the draft's coefficient-dominant order, the gauge
  lies in $\mathcal H$, but $\mathcal H$ is not a formal Laurent series field
  over its constants and Levelt--Turrittin is unavailable. Concretely, the
  $e_z$-component of the Hahn valuation is not a valuation on $\mathcal H$:
  for $f=x^{(0,0)}+x^{(1,-100)}$ and $g=-x^{(0,0)}$ its values on $f$ and $g$
  are $0$ while its value on $f+g$ is $-100$. So there is no Newton polygon, no
  slope filtration and no Turrittin algorithm. Nor can one descend to the
  bounded-$z$-order elements: $f=1-zx^{-\gamma}$ has inverse
  $-\sum_{k\ge1}z^{-k}x^{k\gamma}$, so they do not form a field.
\end{itemize}

An abstract Picard--Vessiot ring for a module over $\mathcal H$ does exist, with
constants $\mathcal C$ after passing to the divisible hull. So the missing
ingredient is \emph{not} a solution algebra. It is a canonical Levelt--Turrittin
normal form over $\mathcal H$, singling out which element of the differential
Galois group deserves to be called the turn. That is the precise shape of the
remaining problem, and it also shows why a purely valuation-theoretic argument
cannot supply it.


\section{The receiver: tensor product, and a better receiver}
\noindent\emph{Archived from the live memo, lines 433--462.}

\subsection*{Why the tensor product cannot simply be dissolved}

The map $a\otimes h\mapsto a\cdot h$ into the Hahn field with $\Omega_{V,j}$
coefficients is not well defined: the copy of $H_{0,j}$ inside $\Omega_{V,j}$
consists of degree-zero constants, the copy inside $H_{\mathrm{tot},j}$ of
monomial series, and the tensor product exists precisely to identify them. Nor
can the two copies be assigned different jobs, with module coefficients read as
constants and gauge bookkeeping as monomials: the gauge relation
$\Omega'=G\Omega G^{-1}-\theta(G)G^{-1}$ ties $G$ and $\Omega$ through the same
Novikov variables, so one copy is forced. With the constants reading the gauge
leaves the field; with the monomial reading the exponential factors become
$\Gamma$-supported and Remark~\textup{[rem:constants-fails]} applies.

Two things should also be said about $\mathscr R_j$ itself. It is a tensor
product of two field extensions and need not be a domain, and the draft reads
characteristic polynomials and algebraic multiplicities in $\mathscr R_j[X]$;
that deserves an explicit remark, which the draft partly gives by proving each
factor injects.

\subsection*{A better receiver}

Extend the $\Gamma$-valuation of $H_{0,j}$ to its algebraic closure and then to
$\Omega_{V,j}$, and take the field of series $\sum_kc_kz^k$ with
$c_k\in\Omega_{V,j}$ whose support $\{(v_\Omega(c_k),k)\}$ is well ordered in
the chosen order. This has a single copy of the coefficients, is a field rather
than a tensor product of two fields, hosts the pro-Laurent gauge, and hosts the
exponential factors as honest units. It is strictly better than (4.4c). It does
not by itself supply Levelt--Turrittin, and it does not remove the rate
criterion below.


\section{The iterated receiver}
\noindent\emph{Archived from the live memo, lines 544--668.}

\section{The iterated receiver}


The criterion above is a consequence of one architectural choice, and a reader
of the second version of this memo proposed the alternative: the draft puts the
intrinsic Novikov variables of the endpoint and the new completion variables
into a single ordered group, which is what gives an eigenvalue difference
positive weight. Nothing forces that. Treat the intrinsic Novikov fraction field
as \emph{coefficients}, and complete only in the new exceptional and bulk
directions afterwards.

Two facts make the proposal concrete, and both are visible in the draft.

\begin{itemize}
\item \emph{The unbounded negative $z$-order is bulk-driven, not
  Novikov-driven.} In the recursion (4.1a) the coefficient $G_\alpha$ has no
  power of $z$ below $-|\alpha|$, where $\alpha$ indexes the \emph{bulk}
  parameter $\eta$, and the filtration depth is acquired by substituting
  $\eta\in F^1B$. A completion in the bulk directions alone therefore already
  hosts the pro-Laurent gauge; the Novikov directions are not needed for it.
\item \emph{With the Novikov directions as coefficients, the pathology
  disappears and the criterion is satisfied rather than evaded.} Eigenvalues of
  $c_1\star$ become constants of weight zero, so $\exp(-\lambda/z)$ is not a
  unit of the receiver but a genuine exponential symbol, the constant-coefficient
  hypothesis of Definition~\textup{[def:solution-algebra]} holds, and
  Remark~\textup{[rem:constants-fails]} evaporates. Since $w(\Delta\lambda)=0$, the
  inequality $e\,w(\Delta\lambda)<\varepsilon$ holds trivially, and the product
  $GP$ of Section~\textup{[sec:receiver]} converges coefficientwise.
\end{itemize}

The shifts that do involve the exceptional Novikov root do not force those
directions back into the filtration. The rank-one unit twist attached to
$-(c-1)\lambda_j$, the fixed divisor shift, and the ambient target's
$\tau^\circ=O(q^{-1})$ all enter Lemma~4.5 as its $a_0$ and $a_2^\circ$ terms,
which are required only to be \emph{constant in} the positive filtration, not to
lie in it. With the Novikov directions as coefficients, $\exp(-a_0/z)$ is an
honest exponential symbol rather than the object the draft must argue about
separately.

\subsection*{What the sources actually permit}

I have checked this against the pinned sources rather than left it as a
question, and the answer is positive for the comparison and negative for the
gauge.

Iritani's Theorem~5.18 states $\Psi$ as an isomorphism of
$\C[z](\!(q^{-1/s})\!)[[Q,\widetilde\tau]]$-modules, with the change of
variables defined over $\C(\!(q^{-1/(r-1)})\!)[[Q]]$. So the ring is
\emph{Laurent} in the exceptional root and \emph{formal} in the ambient Novikov
variables and the bulk. Two consequences.

\emph{The comparison survives.} Inverting the Novikov monomials and completing
in the exceptional and bulk directions is a ring homomorphism out of that ring.
$\Psi$, its inverse, and the two inverse identities are module maps over it, so
they base-change along it unchanged, and the framed operators' entries lie in
the coefficient field, which injects. Nothing in the transport step needs
Iritani's proof, only his statement, so the formality of his ring in $Q$ is not
by itself an obstacle.

\emph{The gauge does not.} The source's own reconstruction writes
$\tau(\widetilde\tau)=\tau^\circ+t(\widetilde\tau)$ and
$\varsigma_j(\widetilde\tau)=\varsigma_j^\circ+s_j(\widetilde\tau)$, with
$t(\widetilde\tau)\in H^*(X)(\!(q^{-1})\!)[[Q,\widetilde\tau]]$ and $s_j$ of the
same shape over $H^*(Z)$. The initial terms are harmless: $\tau^\circ$ and
$\varsigma_j^\circ$ are by definition the values at $Q=\widetilde\tau=0$, hence
carry no Novikov dependence at all, and the exceptional direction remains a
completion direction as the draft already has it. But $t$ and $s_j$ vanish only
at $Q=\widetilde\tau=0$, so at $\widetilde\tau=0$ they retain pure-Novikov
terms. Those have to be removed by the normalized gauge of Lemma~4.5, and that
construction converges by topological nilpotence of the bulk shift. Inverting
the Novikov directions destroys exactly that.

So the iterated receiver is not free, and the decisive question is now small and
graded rather than general:

\begin{quote}
  Is the pure-Novikov part of the coordinate shift --- $t$ and $s_j$ restricted
  to $\widetilde\tau=0$ --- confined to $H^0\oplus H^2$?
\end{quote}

If it is, the string and divisor equations absorb it into the $a_0$ and
$a_2^\circ$ terms of Lemma~4.5, which are required only to be constant in the
positive filtration, and the iterated receiver goes through for the ambient
summand with $w(\Delta\lambda)=0$. If it is not, the Novikov directions must
stay inside the filtration, eigenvalue differences reacquire positive weight,
and the criterion of Section~\textup{[sec:receiver]} returns.

The grading answers this, and the answer is negative in general. Iritani's
Theorem~5.18(5) states that $\tau(\widetilde\tau)$ and
$\varsigma_j(\widetilde\tau)$ are homogeneous of degree $2$, and his conventions
are $\deg Q^d=2c_1(X)\cdot d$ and $\deg\tau_i=2-\deg\phi_i$. A pure-Novikov term
$Q^d\phi$ therefore satisfies
\[
  \deg\phi=2-2\,c_1(X)\cdot d,
\]
so $c_1\cdot d=1$ contributes in $H^0$, $c_1\cdot d=0$ in $H^2$, and
$c_1\cdot d=-1,-2,\dots$ in $H^4,H^6,\dots$; classes with $c_1\cdot d\ge2$
cannot contribute at all. Hence the required confinement to $H^0\oplus H^2$
follows exactly when every contributing effective class satisfies
$c_1\cdot d\ge0$, and it is not automatic: Iritani's theorem is stated for an
arbitrary smooth projective $X$, and after several blowups in a weak
factorization there is no reason for the intermediate ambient variety to have
nef anticanonical class.

So the iterated receiver does not by itself remove arbitrary intermediate
fourfolds from the problem. It removes them under a semipositivity hypothesis on
the ambient, which the weak factorization does not supply. This is worth stating
plainly, because it is the second proposed repair that fails on inspection, and
it fails for a reason the draft would have inherited silently.

For the center the same analysis applies, with the additional and unchanged
difficulty that its Novikov map is the noninjective one, which is why divisor
tagging exists.

The scope consequence is the reason to try this first. If the ambient endpoint
is handled this way, no inequality is needed for arbitrary intermediate
fourfolds --- the case that looked least tractable. What remains is the center
summands, and there the draft's own low-dimensional classification is
restrictive: for nef canonical class it produces an integral-$z$ gauge to a
regular-singular connection, so there is no exponential separation at all and
$w(\Delta\lambda)=0$ again; the remaining minimal cases are $\PP^1$, $\PP^2$,
ruled surfaces, and point blowups. Proving primitive-sixth vanishing directly
for those, in the coordinates Iritani's specialization actually produces, is a
smaller problem than a general transport theorem.


\section{The framing route as proposed, and checks one to three}
\noindent\emph{Archived from the live memo, lines 669--835.}

\section{The framing route, checked against the sources and closed}


The second version of this memo proposed the following route and called it the
probable repair. I state it as it was proposed, then report the three checks,
then state what survives. The checks were run against
arXiv:2411.02266 as cached, arXiv:2307.03696v4, and arXiv:2307.13555v3; the
full verdicts, locators and file hashes are in the companion report
\emph{The three framing-compatibility checks, run against Iritani's
decomposition}.

\subsection*{The route as proposed}

Everything above treats the pro-Laurent bulk gauge of Lemma~4.5 as unavoidable.
It may not be. Hinault, Yu, Zhang and Zhang work with $F$-bundles on the formal
$u$-disc, where a framing is a regular flat connection and, in a framing-flat
trivialization, the meromorphic connection takes the normal form
\[
  \nabla_{u\partial_u}=u\partial_u+u^{-1}K+(\mu D+H).
\]
Their Theorem~4.34 characterizes when two connections of that shape are
gauge-equivalent under
\[
  \Phi(u)\in\mathrm{GL}(H[[u]]),
  \qquad \Phi_{ij}(u)\in k[c_1,\dots,c_r][[u]],
\]
with $\Phi$ then unique given its value at $u=0$; the conditions are that the
leading operators $K$ be conjugate by a parameter-polynomial $\phi$, that the
two $\mu$ agree, and that the diagonal entries of $H$ agree after conjugation
modulo the parameter ideal. Their applications are uniqueness of the
decomposition for projective bundles and for blowups, complementing the
existence theorems of Iritani--Koto and Iritani.

The shape of that gauge is exactly what this memo has been missing. A gauge in
$\mathrm{GL}_n(K[[u]])$ has only nonnegative integral powers of the loop
coordinate, adjoins no root of $u$, and is therefore fixed by the turn
$u\mapsto e^{2\pi i}u$. Taking $\Gamma=0$, so that $\mathcal H=\Omega_V(\!(u)\!)$
and Levelt--Turrittin is available, Theorem~\textup{[thm:transport]} applies with no
further work:
\[
  \Phi\in\mathrm{GL}_n(K[[u]]),\quad
  \nabla'=\Phi\nabla\Phi^{-1}-\theta(\Phi)\Phi^{-1}
  \ \Longrightarrow\
  M_{\mathrm f,V}(\nabla')\ \text{and}\ M_{\mathrm f,V}(\nabla)\ \text{conjugate}.
\]
No Hahn receiver, no product $GP$, and no inequality
$e\,w(\Delta\lambda)<\varepsilon$, provided the decomposition of
Proposition~4.7 can be identified with an $F$-bundle isomorphism of the class
Theorem~4.34 governs. That identification, the matching of the $u=0$ datum, and
the separation of the center specialization were the three checks.

\subsection*{Check one: the identification does not exist as posed}

It fails three times over, and the first two are independent of the draft.

\emph{Domain.} Section~4.4 of Hinault--Yu--Zhang--Zhang is titled
\emph{Equivalence of $F$-bundles over a point}, and both of its applications
evaluate at the closed point $q=t=0$ of the base, the large-radius limit; the
proof of their Theorem~5.16 opens by observing that there the quantum product
reduces to the classical cup product, so that the leading operator $K_{\rm
split}$ of their (5.18) is a cup-product operator, hence scalar-plus-nilpotent
on each block rather than nilpotent outright. The
draft's $\nu_6$ is not evaluated there. It is read where the eigenvalues of
$c_1\star$ separate --- at the cubic, where the block carries $r=(3q)^{1/2}$ and
the indicial roots $\pm1/6$, $\pm5/6$. At $q=0$ that separation collapses and
the exponential blocks the framed operator is built from do not exist. Their
Theorems~5.20 and 5.24 do reconstruct the isomorphism over the whole base from
its restriction at the limit point, but that is uniqueness of a map whose
existence is Iritani's and Iritani--Koto's theorem; it is not a transport of
framed monodromy between base points.

\emph{Hypotheses, for the blowup.} Section~4.4 assumes the generalized
eigenspaces of $K$ all have the same dimension, and the matrix calculus of
Definition~4.33 and Theorem~4.34 rests on the resulting splitting
$H_0^{\oplus m}\cong H$. The blowup decomposition is
$H^*(X)\oplus\bigoplus_{i=1}^{m-1}H^*(Z)[-2i]$, whose summands differ in
dimension whenever $\dim H^*(Z)\ne\dim H^*(X)$; the authors write a general
endomorphism there as a matrix with $\Phi_{1,1}\in\operatorname{End}(H^*(X))$
and $\Phi_{i,i}\in\operatorname{End}(H^*(Z))$, so the mismatch is explicit in
their own notation. Their blowup statement, Theorem~5.22, does assert existence
of the isomorphism, and Theorem~5.24 gives the corresponding uniqueness; the
point is not that it is unproved but that it is not \emph{derived from}
Theorem~4.34, whose hypotheses do not cover the unequal block sizes, and that
its existence is referred to the earlier decomposition work. So Theorem~5.22
supplies no leverage of its own on the blowup half of Proposition~4.7 ---
equation (4.3), the half the headline theorem cannot do without --- beyond what
Iritani's theorem already gives. The projective-bundle half is genuinely covered:
there the eigenspaces are $m$ copies of $H^*(X)$ by their Lemma~5.8, matching
Section~4.4, and Theorem~5.16 is proved.

\emph{Purpose.} Even where Theorem~4.34 applies, its conclusion compares the
source at its base point with the target based at a shifted class $\Delta(a)$,
pinned by their (5.17) for projective bundles and (5.23) for blowups. That
shift is exactly what the pro-Laurent gauge of Lemma~4.5 exists to undo. The
route relocates the bulk displacement rather than removing it.

The two normalizations do agree on where the displacement sits, which is a
useful check on the draft's reading of the sources. In their limit-point
coordinates the eigenvalue $\lambda_i\in\C[c_1,\dots,c_m]$ is cohomology-valued,
so $\Delta(a)$ acquires components in $H^{\ge4}$ from the higher Chern classes,
while the $H^2$ components cancel identically --- their Lemma~5.8(2) gives
$m\lambda_i\equiv m\xi^{i-1}-c_1V$, cancelling the $c_1V$ on the right of
(5.17) --- and the $H^2$ part of $\Delta(a)$ is left free. In Iritani's
coordinates the same data reads
$\varsigma_j^\circ=-(c-1)\lambda_j+h_{Z,j}+O(q^{-1/(c-1)})$ with $\lambda_j$ a
\emph{scalar} multiple of $q^{1/(c-1)}$, so the $H^0$ part is that scalar, the
$H^2$ part is $h_{Z,j}$, and everything of degree at least four has moved into
the tail. Same content, different bookkeeping; the draft's split of
$\varsigma_j^\circ$ into unit twist, fixed divisor, and positive-filtration tail
is right. The tail is present in general, but \emph{not} always: an earlier
version of this passage inferred from the $O(q^{-1/r})$ symbol that even
$X\times\PP^1$ carries a tail and so does not escape the bulk gauge. That
inference was unsupported --- the $O(\cdot)$ symbol carries no information about
$[z^{-1}]\log$ either way --- and it is contradicted by
Theorem~\textup{[thm:trivial-displacement]} below, which shows that for a trivial
bundle $\varsigma_j^\circ=r\lambda_j$ \emph{exactly}, with no $H^2$ term and no
$H^{\ge4}$ tail. The claim is struck.

\subsection*{Check two: the $u=0$ datum, which passes}

Their Theorems~5.20(2) and 5.24(2) determine the base map from its restriction
to the base point up to a multiplicative constant in the $q$ direction (the
statement for all logarithmic directions is their Proposition~4.31(2));
the proof of their Proposition~4.31 shows why, since $df$ is recovered through
$\Psi(d\log p_i)=d\log f_i$ and integrating a logarithmic form leaves one
constant. That ambiguity is a character of the Novikov lattice: $q\mapsto cq$ is
$Q^d\mapsto c^{E\cdot d}Q^d$, which is the divisor substitution (4.1) of the
draft with $a_2^\circ$ a logarithmic multiple of the dual class. The draft
already proves such a substitution is a $z$-constant coefficient automorphism
and leaves $\nu_6$ unchanged, and already checks it is well defined on the image
monoid. The bundle map carries no ambiguity at all once its restriction is
fixed.

The ambiguity also never has to be tolerated, because the sources pin it.
Iritani--Koto's (5.11) gives
$\varsigma_j^\circ=r\lambda_j+[z^{-1}]\log\bigl(q^{c_1(V)/(rz)}F_j(1)\bigr)$ in
closed form and their (5.12) gives $\Phi_j^\circ$; Iritani's Section~5.8.1 fixes
$\Psi|_{Q=\widetilde\tau=0}$ from his Lemma~5.13 and the explicit $\lambda_j$.
One manuscript sentence citing those, and noting that any residual constant is a
divisor substitution already covered by (4.1), settles this point for good,
whichever repair eventually lands.

\subsection*{Check three: the center specialization, which stays separate}

It stays separate more firmly than the check anticipated. Their $A$-model
$F$-bundle is built throughout on the one-variable Novikov projection
$q^\beta\mapsto q^{\beta\cdot\omega}$ of their (2.21), and in the blowup
application the center copies are the maximal $F$-bundles of $Z$ associated to
$\sigma^*\omega_X$ --- the Novikov direction pulled back from the ambient
variety, which is the draft's specialization in a coarser form. No statement of
theirs can therefore certify that a specialized center module carries an
intrinsic framed monodromy. Divisor tagging remains the route, in the current
manuscript, for identifying the specialized center contribution with the
intrinsic invariant; route~3 of the list below bypasses that identification
instead, by proving specialized vanishing directly in low dimension.

One genuine gain came out of it. Their Lemma~2.24 proves that the collapsed
potential determines the full genus-zero potential, by expanding in the $H^2$
bulk directions and separating curve classes by their intersection numbers with
an $H^2$ basis, under the finiteness of their Assumption~2.22. That is the
mechanism of \texttt{lem:divisor-tagging}, in print, in a closely related
setting. It does not close that lemma --- it works at the level of the potential
over a base retaining all $H^2$ directions, whereas the draft needs the
separation at the level of a framed operator at a fixed bulk point, and the
draft's collapse is by the pair $(i_*,\rho_C\cdot)$ rather than by a single nef
class --- but it makes the shape of the argument a cited one.


\section{The corrected program (superseded by the rigidity section)}
\noindent\emph{Archived from the live memo, lines 1121--1150.}

\subsection*{The corrected program}

\begin{enumerate}
\item Flatness preserves the spectral decomposition of $U$ and gives the exact
  evolution $D_aU_i=C_{a,i}+[C_{a,i},\mu_i]$, unconditionally.
\item Multiplicity-one blocks contribute no cyclotomic monodromy
  (Theorem~\textup{[thm:simple-blocks]}), so only coalesced blocks are in play.
\item For a coalesced block the off-block projector equations are governed by the
  full Sylvester operator $(u_j-u_i)+L_{N_j}-R_{N_i}$, invertible but not scalar.
\item What must be deformed is the entire formal block normal form
  $J+zD+z^2E+\cdots$, not the compressed pair, and the deformation theory has to
  reproduce the cubic $2\times2$ example above.
\end{enumerate}

A caution on the literature before it is leaned on. The extension results for
isomonodromic deformations with coalescing eigenvalues assume, among other
hypotheses, that the leading matrix is diagonalizable at the coalescence point.
The carrier here is nonsemisimple at leading order --- $J_0\ne0$ --- so those
theorems do not apply off the shelf, and the coalescence question stated below
is not answered by citing them.

The risk that remains is unchanged in kind and now better located: a coalesced
block can \emph{split} under the displacement, the splitting being governed by
$D_aU_i$, and if the rank-two block carrying $\pm1/6$ split into simple blocks
those would carry residue $0$ by Theorem~\textup{[thm:simple-blocks]} and $\nu_6$
would drop by two. Proving that the reconstruction displacement does not do
that, or does it in a count-preserving way, is the residual problem.
Section~\textup{[sec:rigid]} now does exactly that, by shearing the block before
deforming it.


\section{A shorter route to the endpoint via the Serre decoration (not taken)}
\noindent\emph{Archived from the live memo, lines 1542--1740.}

\subsection*{A shorter route to the endpoint, in the sources' own terms}


Two paragraphs of Katzarkov--Kontsevich--Pantev--Yu's Example 6.21 bear directly
on this and should be on the record, because between them they suggest a route
to the one-stabilization theorem that needs neither this memo's transport lemma
nor its rigidity theorem.

First, they compute the atomic composition of a smooth cubic threefold at the
hyperplane point: two one-dimensional atoms for the nonzero eigenvalues of
$\operatorname{Eu}\star$, and one atom $\alpha(X)$ for the eigenvalue $0$. That
is the block structure used throughout this memo, in their language. They add
that ``the Witt algebra argument from Remark~3.14 shows that there is no further
splitting in a neighborhood of $b$'' --- which is
Theorem~\textup{[thm:no-splitting]}, asserted rather than proved, and their
Remark~3.14 introduces its finer clauses with ``more fine results can be
proved''. Section~\textup{[sec:rigid]} supplies proofs of both the splitting
statement and the finer one, for this block type, from flatness alone.

Second, and more consequentially, they separate $\alpha(X)$ from a curve atom
not by formal monodromy but by a Serre automorphism. They observe that
$\alpha(X)$ looks like the atom of a smooth projective genus-five curve, then
note that as a Hodge atom enhanced with its Serre automorphism $S$, the graded
minimal polynomial forbids $S$ from having eigenvalue $-1$, so it cannot be the
Serre automorphism of a genus-five curve. They conclude that every smooth cubic
threefold is irrational, and remark that the same approach applies to other Fano
threefolds.

Their $S^5=[3]$ is quoted correctly above --- it is what Example 6.21 says --- but
it is not the relation to use when computing on a K-group. Kuznetsov's fractional
Calabi--Yau relation for the same category is $S^3=[5]$, and on the numerical
Grothendieck group, where a shift $[k]$ acts by $(-1)^k$, the two predict
$S^3=-I$ and $S^5=-I$ respectively; the Serre operator computed from
Riemann--Roch satisfies the first exactly and fails the second, its
characteristic polynomial being $\Phi_6$ and its order six, in agreement with the
exponents $-1/6,-5/6$ recorded after \textup{[eq:sheared-data]}. The offset is
systematic rather than a slip: their Example~6.19 gives $S^3=[4]$ for a very
general quartic in $\PP^5$, whose Kuznetsov component is a K3 category with
$S=[2]$ categorically. (An earlier version of this passage cited an
``Example~6.20'' about a cubic fourfold; no such example exists, and the
statement belongs to Example~6.19 and to the quartic.)
The two are different objects --- the Serre automorphism of a Hodge atom in the
cohomological $\mathbf Z$-grading, and the categorical Serre functor on a K-group
--- so a citation should follow the object it is about. Locators and the full
comparison are in `2026-08-15-c912-kkpy-imports-source-check.md'.

A second normalization difference runs alongside it and explains a substitution
this lane had observed but not accounted for. Their statements are made after the
half-parity gauge $u^g$ of Claim~6.15, this memo's count before it. That gauge
shifts exponents by half the cohomological degree, multiplying monodromy
eigenvalues by a parity-dependent sign: primitive sixth roots become primitive
cube roots and $\pm1$ swap. It is the $\lambda\mapsto-\lambda$ substitution
between the prime-Fano classification's polynomial $R$ and the Serre side. Every
separation used below distinguishes roots of unity of order at most two from
roots of order three or six, which the gauge preserves.

The consequence for the endpoint is worth stating plainly. Their argument is
dimension three, where Proposition~5.30 requires excluding only points and
curves. The one-stabilization statement is dimension four, where the same
proposition requires excluding points, curves \emph{and surfaces}. Everything
else --- the projective-bundle formula putting two copies of $\alpha(X)$ in
$X\times\PP^1$, the non-rationality criterion, and the decoration itself ---
is already theirs. So the endpoint theorem may reduce to one classification
question, which in the terms of this memo reads:

\begin{quote}
\textbf{(iii)} Every atom of a smooth projective surface has formal monodromy of
order at most two, its eigenvalues lying in $\{+1,-1\}$. In particular no such
atom has a primitive sixth root of unity as an eigenvalue, and no Serre operator
on such an atom's numerical Grothendieck group has $\Phi_6$ as a factor.
\end{quote}

The stronger form is the one Katzarkov--Kontsevich--Pantev--Yu's own dimension-four
example uses, where the exclusion is unipotency against non-unipotency. Their
Example~6.19 needs it only for surfaces of general type, the atom there being
shaped like one; the endpoint needs every surface, because the cubic threefold's
atom is shaped like a genus-five curve atom and nothing narrows the search in
advance.

Two reductions come with the statement, and one warning.

\emph{Rank-one atoms are automatic.} For a rank-one lattice the Euler form is a
$1\times1$ matrix, hence symmetric, so $S=E^{-1}E^T=1$. A semisimple small
quantum cohomology has every block of rank one, so (iii) holds for every surface
with semisimple small quantum cohomology with nothing to prove. For a surface
with nef canonical class their Lemma~5.24 gives a single atom and Claim~6.15
gives a regular singularity with nilpotent residue after the half-parity gauge,
putting every eigenvalue at a root of unity of order at most two.

\emph{The remaining surfaces, by the minimal model program.} These two
reductions leave a short induction, which closes (iii) modulo three imported
statements.

\begin{remark}[This is Theorem~\textup{[thm:low-dim]}, and the two readings coincide]

Statement (iii) is not new work: it is Theorem~\textup{[thm:low-dim]}, proved in
Section~\textup{[sec:endpoint]} by exactly this induction --- blow down to a minimal
model, use their Lemma~5.24 and Claim~6.15 for nef canonical class, and reduce
$\PP^2$ and the geometrically ruled surfaces to point and curve atoms through the
projective-bundle identity. What the present section adds is only that the
Serre-decorated reading of (iii) is the same statement rather than a companion to
it: Katzarkov--Kontsevich--Pantev--Yu \emph{define} the Serre automorphism as the
monodromy in the $u$-direction, with $\chi(a,b)=\chi(b,S(a))$, so an atom's Serre
eigenvalues are its formal monodromy eigenvalues and the shorter route's step
(iii) is discharged by Theorem~\textup{[thm:low-dim]} as it stands.

Two refinements are worth recording. First, the conclusion is sharper than
$\nu_6=0$: in every case the proof produces monodromy of order at most two, which
is the discriminator their own dimension-four example uses. Second, the blowup
input can now be taken from Iritani's Theorem 5.18 directly, an isomorphism of
quantum $D$-modules commuting with the connection and preserving the Euler
vector fields, rather than only through the chemical-formula identity.
\end{remark}

The fake of the previous paragraph does not disturb the blow-down step: it mixes
the new exceptional object with the ambient structure sheaf, which lie in
different atoms, so the pair it generates is not an atom. Nor is the blow-down
step resting on a multiplicity count: the computational literature --- Böhning,
von Bothmer and Su'a --- verifies the blowup statement for multiplicity
sequences and records that naive and small atomic decompositions can differ, but
Iritani's theorem is at the level of the connection itself, which is the reading
used above. The reading list, cache locators and hashes are in
`2026-08-15-c912-blowup-formula-source-check.md'. Iritani also records that
Gyenge and Szab\'o studied the asymptotics of the Euler eigenvalues under
blowdowns to minimal models of surfaces, which is the configuration this
induction runs through.

\emph{The ambient category never sees it.} The Serre functor of a surface is
tensor by $\omega$ followed by $[2]$, and $[2]$ acts by $+1$ on K-groups, so the
ambient numerical Serre operator is multiplication by $\exp(K)$, which is
unipotent. A primitive sixth eigenvalue on a surface can only come from a
decomposition, never from the surface's own K-theory.

\emph{The warning: admissible is not enough.} For a subcategory generated by an
exceptional pair with $\chi(e_1,e_2)=x$ the numerical Serre operator has
characteristic polynomial $\lambda^2+(x^2-2)\lambda+1$, which is $\Phi_6$ exactly
when $x^2=1$. On any surface carrying a $(-1)$-curve $E$ --- that is, any
non-minimal surface --- the sheaf $O_E(-1)$ is exceptional,
$\operatorname{Hom}^\bullet(O_X,O_E(-1))=0$, and $\chi(O_E(-1),O_X)=-1$, using
only $E^2=-1$ and $K_X\cdot E=-1$.
So $\langle O_E(-1),O_X\rangle$ generates an admissible subcategory whose Serre
operator has order six. Statement (iii) is therefore false if ``atom'' is
weakened to ``admissible subcategory'', and no lattice argument ranging over
admissible subcategories can prove it. The subcategory is not an atom: it is
generated by an exceptional pair and splits into two rank-one pieces whose Serre
operators are both the identity. Since the blowup is exactly what supplies this
fake, the blowup step of any proof of (iii) is where the care is needed.

With all four imports now read at the source, the route to
``$X\times\PP^1$ is not rational'' runs in four steps. Their Example~6.21 gives
the atom $\alpha(X)$ at the eigenvalue zero. Its formal monodromy has the two
primitive sixth roots of unity as eigenvalues, by the exponents $-1/6,-5/6$
computed after \textup{[eq:sheared-data]}; that is order six. Theorem~4.11 with
$r=2$, applied to $X\times\PP^1=\PP(\mathcal O\oplus\mathcal O)$, puts two copies
of $\alpha(X)$ in the atomic decomposition of $X\times\PP^1$. And no atom of a
smooth projective variety of dimension at most two has monodromy of order six:
points give the identity, curves give $-1$ twice, and surfaces have order at most
two by Theorem~\textup{[thm:low-dim]}.

The criterion applied at the end is \emph{not} their numbered Proposition~5.30,
which excludes membership in $\mathrm{HAtoms}$ of low dimension for undecorated
atoms, whereas the separation just made uses the monodromy, a decoration. It is
the enhanced form, which the same section states as unnumbered prose: if a smooth
projective variety has an enhanced atomic composition containing an enhanced atom
that does not come from a smooth projective variety of dimension at most
$\dim X-2$, then it is not rational. With $\dim(X\times\PP^1)=4$ the excluded
dimensions are at most two, which is what the four steps supply. The numbered
proposition is proved there as a special case of their Proposition~5.17; the
enhanced version is asserted, on the ground that enhancement by Euler pairings
and Serre automorphisms repeats the undecorated theory. So the load-bearing
citation of this route is an assertion, and a manuscript using it must either
prove the enhanced criterion in the generality it needs or say plainly that it
quotes one. The route uses no transport lemma, none of
Section~\textup{[sec:unconditional]}, and no integral structures.

This is not the same invariant as $\nu_6$, and it is not obviously weaker or
stronger; it is the sources' own, which is an argument for using it. Two
cautions. It buys the Hodge-atom, motive and Serre-enhancement machinery that
the earlier blueprint deliberately avoided, so the manuscript would import more
than the two decomposition theorems. The deferral is wider than an earlier
version of this passage recorded: in the non-archimedean analytic setting, which
is the one an $A$-model $F$-bundle of a projective variety actually lives in, the
authors defer the Euler pairing and the Serre automorphism themselves to
forthcoming work, not merely the integral-structure enhancement. So the
decoration the route turns on is not available off the shelf in the setting
required. And their sentence is a worked example,
not a theorem with hypotheses, although the statements it invokes ---
Lemma~5.24, Claim~6.15, Proposition~5.30 and Theorem~4.11 --- have now been read
at the source and say what is claimed of them here. The graded minimal
polynomial has been checked on the lattice as well, and the two conventions
differ as described above; the surface exclusion is
Theorem~\textup{[thm:low-dim]}, with Remark~\textup{[rem:step-iii-is-low-dim]} for why the
decorated reading needs nothing further.
Backing computations, and the exact record of what was and was not read at the
source, are in `2026-08-15-c912-step-iii-serre-restatement.md' and
`2026-08-15-c912-step-iii-surface-induction.md', with the scripts
`2026-08-15-c912-surface-atom-serre-check.py' and
`2026-08-15-c912-step-iii-surface-induction-check.py'.


\section{Which other routes also close (merged into the routes section)}
\noindent\emph{Archived from the live memo, lines 2177--2259.}

\subsection*{Which other routes also close}

Three, at different costs.

\emph{The draft's own presentation.} Keep the operation formulas and transport
$\nu_6$ summand by summand. This needs the specialization-safe statement at
every intermediate fourfold, so it needs Section~\textup{[sec:rigid]} generalized
from a rank-two Jordan block to an arbitrary block. That generalization is
plausible for a single Jordan block of any size by the same shearing, and needs
a resonance hypothesis in the semisimple case; it is real work, and it is the
only route that keeps the draft's current architecture.

\emph{The sources' Serre decoration: restricted, referee-risky, and not taken.}
On its face it avoids the gap of Section~\textup{[sec:endpoint-gap]} by construction,
which is what made it look attractive; the reasons not to take it are recorded
after the description. Their Proposition 5.23 proves that the
isomorphism class of the fibre representation of an atom is independent of the
representative, using rigidity of finite-dimensional representations of a
proreductive group --- a global argument, valid across a connected component,
which monodromy data does not enjoy. They then decorate with a Serre
automorphism and state that the enhanced theory is a straightforward repeat of
the undecorated one, with the corresponding strengthened non-rationality
criterion.

Their Example 6.21 separates $\alpha(X)$ from a genus-five curve atom by the
graded minimal polynomial $S^5=[3]$ and concludes that every smooth cubic
threefold is irrational. More to the point for the endpoint, their Example 6.17
runs the same pattern in \emph{dimension four}, for a very general cubic
fourfold containing a plane: there $\alpha(X)$ is shaped like the atom of a
surface of general type, and they exclude it by playing $S^3=[4]$ against the
unipotency that Claim 6.15 forces on the even cohomology of a nef-canonical
surface. That is exactly the shape the one-stabilization argument needs, with
$X\times\PP^1$ in place of their fourfold, and it pairs their globally
well-defined decoration on the cubic side with a pointwise structural statement
on the surface side --- the two halves whose mismatch is the gap above.

Why it is not taken, in order of severity.

\emph{The decoration is not defined in the category our atoms live in.} They
define the duality automorphism $S$ for a \emph{complex analytic} F-bundle of
exponential type with nondegenerate pairing, as the monodromy in the $u$
direction, and then say that ``the definition in the case of non-archimedean
analytic F-bundle requires extra work, and will be discussed in [49]''. But the
A-model F-bundle of a projective variety is built over $k((y))$ and base changed
to its algebraic closure, and is a non-archimedean analytic F-bundle in their own
words. So in the setting where Hodge atoms of a projective variety live, the
decoration this route turns on is not merely unproved but undefined as the paper
stands. That is not a gap a lemma of ours can close; supplying the definition is
the content of their deferred work.

\emph{The enhanced theory is asserted, not written out.} The sentence that the
enhanced theory ``is completely straightforward and is just a repeat of the
theory of undecorated Hodge atoms'' is unnumbered running text. The
corresponding strengthened non-rationality criterion, which is the one this route
would apply, is stated in the same register; the numbered and proved
Proposition~5.30 is the undecorated version, which does not separate atoms by a
decoration. A manuscript may not stand on the unnumbered version.

\emph{The statements live in three different categories.} Their Leray--Hirsch
Theorem~4.11 is an isomorphism over analytic domains, their atomic statements
hold near a chosen base point, and Iritani's blowup decomposition is formal over
an extended Novikov ring with fractional powers of $q$. Any argument of ours
would stitch formal, non-archimedean analytic and complex analytic statements
together, and every stitch is a place a referee asks precisely what was
transported.

\emph{And the examples are examples.} Their Examples 6.17--6.21 are worked cases,
not theorems with hypotheses, so the surface exclusion for $X\times\PP^1$ would
have to be done rather than cited in any case.

One way the first objection could be dodged is to work complex analytically,
where $S$ is defined, which requires convergence of the relevant quantum
cohomology and so trades Hypothesis 4.7H for a convergence hypothesis. That is a
hypothesis of a more familiar kind, but it is still a hypothesis, and the formal
treatment of this memo was chosen precisely to avoid needing one. The route is
therefore recorded as a restricted and referee-risky alternative with heavy
import costs, not as the plan.

\emph{Both decorations at once.} They are independent invariants of the same
atom, so proving the surface exclusion twice is a genuine referee-facing check
rather than duplication. If the Serre computation is cheap it is worth having as
a second column in the same table.


\section{The trade between the two sides, and where it cancels}
\noindent\emph{Archived from the live memo, lines 2372--2428.}

\subsection*{The trade between the two sides, and where it cancels}


It is worth being explicit about what is being traded, because the accounting
localizes the remaining debt more sharply than the narrative does.

\emph{Vertical motion is free; horizontal motion is the whole debt.} Each step of
a factorization gives an exact identity
$\nu_6(\mathrm{Bl}_ZV,w)=\nu_6(V,\tau(w))+\sum_j\nu_6(Z,\varsigma_j(w))$, and with
$\dim Z\le2$ the centre terms vanish. So moving between varieties conserves the
invariant exactly: nothing leaks into a centre and nothing is created. The only
way the invariant can change anywhere in the ledger is by moving between two
parameters of \emph{one} variety. Every obligation this memo has chased is a
horizontal one, and the vertical steps have never been in doubt.

\emph{Consequence, in contrapositive form.} Suppose $X\times\PP^1$ is rational.
The left end carries $4$ and the right end carries $0$, and conservation across
every vertical step then forces some valley of the zigzag to have a
bulk-parameter-dependent primitive-sixth multiplicity. So the theorem is
equivalent to: no smooth projective fourfold arising as a valley has a
$\nu_6$-carrying block that degenerates along the displacement the chain
produces. Stated this way the debt is a single negative statement about
degeneration, not a positive statement about transport.

\emph{The balance has slack, and the slack is three.} The left end supplies $4$,
being two copies of the cubic's rank-two block, and the contradiction needs only
$1$. So up to three of the four eigenvalues may be lost anywhere along the chain
without harming the argument. What must be excluded is total collapse, not any
loss --- a weaker statement than constancy, and one that Section~\textup{[sec:rigid]}
overshoots, since it proves the characteristic polynomial exactly constant.

\emph{Four places where the debt cancels outright.}
\begin{enumerate}
\item Any step whose accumulated displacement lies in $H^0\oplus H^2$. The string
  and divisor equations are exact there, so such a step costs nothing. This is
  not a rare case: it is exactly what makes the endpoint free, where the
  displacement is $\pm2q^{1/2}$ and purely $H^0$.
\item An adjacent blowup and blowdown of the same centre, where the composite
  parameter map is the identity and no comparison of two parameters arises.
\item Any valley whose Euler operator is semisimple, or more generally has only
  multiplicity-one eigenvalues in the relevant range. There
  Theorem~\textup{[thm:simple-blocks]} gives $\nu_6=0$ at every parameter
  unconditionally, so the two incoming parameters agree trivially and the valley
  is vacuous.
\item Every centre, by the low-dimensional lemma. Centres are where a leak would
  have to go, and they cannot receive one.
\end{enumerate}

\emph{What is left after the cancellations.} The debt sits only at valleys that
simultaneously carry a nonzero $\nu_6$ and receive a displacement with a nonzero
$H^{\ge4}$ part. By the dimension axiom each such displacement contributes
finitely much per Novikov degree, so the surviving obligation is finite data at
finitely many places, and by the slack it need only be shown not to annihilate
all four eigenvalues at once. That is the sharpest available statement of what
remains, and it is what a next attempt should target rather than constancy in
general.


\section{Corrections to the earlier versions of this memo}
\noindent\emph{Archived from the live memo, lines 2570--2712.}

\section{Corrections to the earlier versions of this memo}


\subsection*{From the fifth version to this one}

Nothing in the fifth version is retracted; one of its hypotheses is discharged
and one of its expectations is corrected.

\begin{enumerate}
\item Standing hypothesis (H2) of Section~\textup{[sec:rigid]} is deleted. It is a
  theorem: Lemma~\textup{[lem:duality-gauge]} shows the decoupling gauge is an
  isometry of the Poincar\'e pairing, and Theorem~\textup{[thm:h2-automatic]} then
  forces $f\equiv0$ pointwise on the germ, for every rank-two coalesced block of
  every smooth projective target. Theorem~\textup{[thm:no-irregularity]} and
  Corollary~\textup{[cor:cubic-closed]} now hold under (H1) alone. The flatness route
  to $f\equiv0$ is kept as a remark, since it is what survives without the
  Frobenius property.
\item The shorter-endpoint-route subsection recorded the graded minimal
  polynomial as $S^5=[3]$. On the numerical Grothendieck group that is
  incompatible with the count: it forces order ten, where the Serre operator
  computed from Riemann--Roch has characteristic polynomial $\Phi_6$ and order
  six, matching Kuznetsov's $S^3=[5]$. The subsection now uses the verified
  relation, and states the surface obligation as (iii) with its two reductions
  and the non-minimal-surface counterexample to the weakened form. It also now
  records that the shorter route's step (iii) is discharged by the existing
  Theorem~\textup{[thm:low-dim]}, since their Serre automorphism is defined as the
  monodromy and the two readings therefore coincide
  (Remark~\textup{[rem:step-iii-is-low-dim]}); an intermediate version of this memo
  carried a separate proposition duplicating that theorem, now removed.
\item The scope remark expected the rank-two mechanism to extend to a Jordan
  block of any size. It half does. Proposition~\textup{[prop:size-m-steps]} carries
  the commutant and no-splitting steps to size $m$, but
  Proposition~\textup{[prop:size-m-duality]} shows duality removes only the deepest
  of the pole coefficients the general shearing produces, so rank two is special
  rather than typical, and for $m\ge3$ the base point must be settled by the
  grading, whose missing step is named there.
\end{enumerate}

\subsection*{From the fourth version to the fifth}

Nothing in the fourth version is retracted. Section~\textup{[sec:rigid]} is added,
and three of its consequences change earlier readings.

\begin{enumerate}
\item Section~\textup{[sec:not-yet]} said the coalesced case was not reduced and
  named the block-splitting risk as the residual problem. Both are now settled
  for a rank-two Jordan block: Theorem~\textup{[thm:no-splitting]} excludes the
  splitting and Theorem~\textup{[thm:rigidity]} freezes the exponents. The section is
  left standing because its diagnosis --- that the compressed pair is the wrong
  object --- is what the shearing acts on.
\item The residual gap of Section~\textup{[sec:hyzz]} is closed for the cubic
  ambient summand by Corollary~\textup{[cor:cubic-closed]}. It remains open as
  stated for arbitrary smooth projective $T$, which is a statement about
  arbitrary blocks, not about transport.
\item Cai's Proposition~6 already asserts the $\pm1/6$ exponents for the big
  quantum connection, by a bulk gauge of exactly the pro-Laurent shape this memo
  scrutinized. That argument is sound with the bulk kept formal; the receiver
  discussion was always about specialization, not about the formal statement.
  Section~\textup{[sec:rigid]} makes the point moot by producing constants.
\end{enumerate}

\subsection*{From the third version to the fourth}

Both are in Section~\textup{[sec:blockwise]}, and both are at the coalesced blocks,
which are the ones that matter.

\begin{enumerate}
\item The off-block projector equation was derived as
  $(\partial_aU)_{ji}+(u_j-u_i+N_j)X_{ji}=0$, dropping the term $-X_{ji}N_i$.
  The correct equation is $\mathcal S_{ji}(X_{ji})=-(\partial_aU)_{ji}$ with
  $\mathcal S_{ji}(X)=(u_j-u_i)X+N_jX-XN_i$. The closed form
  $D_a\mu_i=[C_{a,i},S_i]$ therefore holds only when the block itself is
  semisimple, and the blocks carrying $\nu_6$ are exactly the ones where it is
  not; in the draft's cubic block $J_0\ne0$.
\item The block was modelled as $z\partial_z-z^{-1}U_i+\mu_i$. After formal
  decoupling a nonsemisimple block carries higher powers of $z$, and its
  exponents depend on them: in the cubic block the $z^2$ coefficient's entry
  $-8/81$ enters $\rho^2+\rho+5/36=0$, and deleting it moves the exponents from
  $\pm1/6$ to $\pm1/18$. So the isomonodromy criterion stated there is a
  statement about semisimple blocks and is not a reduction of the residual
  problem. It is restated with that hypothesis.
\end{enumerate}

The simple-block theorem survives, with one wording repair: the formal solution
of a simple block is $e^{-u_i/z}$ times a single-valued formal unit, not times a
constant, since decoupling can generate higher $z$-terms there too. The
monodromy conclusion is unaffected.

\subsection*{From the second version to the third}

Both changes come from running the three checks against the sources.

\begin{enumerate}
\item The framing route of Section~\textup{[sec:hyzz]} was called the probable
  repair. It is not a repair at all: its governing theorem lives at the
  large-radius limit point; its blowup case, though stated with existence in
  Theorem~5.22 and uniqueness in Theorem~5.24, is not derived from
  Theorem~4.34, whose hypotheses do not cover the unequal block sizes, and its
  existence is referred to the earlier decomposition work; and the gauge it
  supplies lands at the shifted base point the bulk gauge exists to undo.
\item The second version wrote that the unbounded negative loop order is the
  draft's own choice rather than something the comparison theorems impose. That
  is wrong, and it was the ground for the recommendation. Iritani--Koto
  transport along the same object, their fundamental solution $M$ of
  Section~5.8, with the same per-bulk-degree $z^{-1}$-polynomiality the draft
  derives in (4.1a).
\end{enumerate}

Two findings were gained in exchange, and they narrow the problem rather than
restate it: the comparison isomorphisms need no receiver, and the residual gap
is one displacement statement. Both are in Section~\textup{[sec:hyzz]}.

\subsection*{From the first version to the second}

For readers of the first version, these are the substantive changes.

\begin{enumerate}
\item The Constants Lemma was stated without the constant-coefficient
  hypothesis and is false without it; see Remark~\textup{[rem:constants-fails]}. The
  framed-operator characterization inherits the hypothesis.
\item The scissors were only implicit. The transport theorem has no verified
  instance in the application, and that now appears where the theorem does.
\item The claim that the criterion is scale-invariant, and the inference that
  the separating parameter $L$ cannot be tuned, were both wrong. Tuning $L$ down
  is now the first proposed route.
\item The cubic-endpoint bullet compared two different normalizations. Corrected
  above.
\item The finite-tensor objection to the receiver was wrong twice over: the
  pro-Laurent gauge lies in the Hahn factor alone, and the splitting gauge's
  coefficients lie in one finite extension of $H_{0,j}$. It is replaced by the
  gauge-relation argument.
\item The remark that a tensor product of two fields is not a domain was
  asserted as a general principle, which is false. The correct statement is that
  $\mathscr R_j$ need not be one.
\item The obstruction is now given in its order-free form, which answers two of
  the questions the first version left open: no completion helps, and the
  criterion is not an artifact of the ordered receiver.
\item Smaller repairs: the freeness proof (the relevant quotient group has
  torsion), the commutation of $\sigma$ with the derivation, the inverse of $M$,
  the range of the exponents, the ramification factor in the splitting rate, and
  the divisor-tagging observation, which is new.
\end{enumerate}


\end{document}